% ============================================================================
% Section 1: Introduction
% ============================================================================

\section{Introduction}
\label{sec:introduction}

Consider an enterprise deploying an AI agent for customer support ticket
routing. On Monday, after a prompt refinement, the agent correctly routes
93\% of tickets. By Wednesday, a model provider silently updates the
underlying LLM, and the routing accuracy drops to 71\%. No test caught
the regression. No alert fired. The team discovers the degradation only
after a surge in customer complaints.

This scenario---\emph{works Monday, fails Wednesday}---is not
hypothetical. It is the daily reality for engineering teams deploying
autonomous AI agents in production. The root cause is
\emph{non-determinism}: the same agent configuration (prompt, tools,
model, orchestration logic) can produce different outputs across
invocations due to temperature sampling, model weight updates, tool
latency variations, and context window effects. Recent empirical work
confirms that agent behavior exhibits stochastic path deviation
\citep{kapoor2024capable}, yet no principled testing methodology exists
to detect when this deviation constitutes a \emph{regression}---a
statistically significant degradation in agent quality.

\subsection{The Failure of Traditional Testing}
\label{sec:intro:failure}

Traditional software testing rests on two assumptions that are
fundamentally violated by AI agents:

\begin{enumerate}[leftmargin=*]
  \item \textbf{Determinism.} Given the same input, a function produces
    the same output. Tests assert equality: \texttt{assert f(x) == y}.
    For agents, the same input can yield different tool selections,
    different reasoning chains, and different final answers across runs.
    Equality assertions are meaningless.

  \item \textbf{Binary verdicts.} A test either passes or fails. For
    stochastic systems, the question is not ``did it pass?'' but ``does
    it pass with sufficient probability?'' A single failure may be
    statistical noise; a single success may be a lucky run. The entire
    notion of pass/fail must be reconceived.
\end{enumerate}

Existing approaches to LLM evaluation---deepeval
\citep{deepeval2024}, promptfoo \citep{promptfoo2024}, OpenAI Evals---address
single-turn output quality but do not formalize regression detection
for multi-step agent workflows. They evaluate \emph{how good} an
agent is; they do not verify \emph{whether it got worse}. The recently
introduced agentrial framework \citep{agentrial2026} takes an important
first step by running agents multiple times and computing confidence
intervals, but it lacks a formal theoretical foundation: no stochastic
test semantics, no coverage metrics, no mutation testing, no composition
theory, and no proofs of statistical guarantees.

\subsection{The Paradigm Shift: Binary to Probabilistic}
\label{sec:intro:paradigm}

We argue that agent testing requires a fundamental paradigm shift from
\emph{deterministic, binary} verdicts to \emph{stochastic, three-valued}
verdicts:

\begin{center}
\begin{tabular}{lcc}
\toprule
\textbf{Aspect} & \textbf{Traditional Testing} & \textbf{Agent Testing} \\
\midrule
Execution model & Deterministic & Stochastic \\
Verdict space & $\{\PASS, \FAIL\}$ & $\{\PASS, \FAIL, \INCONC\}$ \\
Assertion & $f(x) = y$ & $\Pr[f(x) \models \phi] \geq \theta$ \\
Regression & $f'(x) \neq f(x)$ & $p_{\text{new}} < p_{\text{base}} - \delta$ \\
Evidence & 1 run & $n$ runs with confidence \\
\bottomrule
\end{tabular}
\end{center}

The third verdict value, \INCONC{} (\emph{inconclusive}), is essential:
when the number of trial runs is insufficient to distinguish signal from
noise, the honest answer is neither pass nor fail but ``more evidence is
needed.'' This three-valued semantics is the foundation upon which all
other contributions in this paper are built.

\subsection{The Cost of Testing Non-Determinism}
\label{sec:intro:cost}

The statistical rigor of stochastic testing comes at a price: each
trial requires a full agent execution, consuming API tokens at
non-trivial cost. To detect a regression of magnitude $\delta = 0.10$
at $\alpha = 0.05$ and $\beta = 0.10$, the fixed-sample formula
(\cref{eq:sample-size}) requires approximately $n \approx 100$ trials
per scenario. A test suite of 50 scenarios at 100 trials each
generates 5{,}000 agent invocations per regression check. For
frontier models where each complex agent run costs \$5--15 in API
tokens, a single regression check costs \$25{,}000--75{,}000---more
than many teams spend on production usage in a month.

This \emph{cost barrier} fundamentally limits adoption. In CI/CD
pipelines where code is deployed multiple times per day, the testing
budget must be orders of magnitude lower than production cost. Even
with SPRT reducing trials by $\sim$50\% (\cref{sec:stochastic:sprt}),
the cost remains prohibitive. No existing framework addresses the
\emph{token economics} of agent testing: how to achieve the same
statistical guarantees at radically lower cost.

We observe that the cost problem arises from a mismatch between the
dimensionality of agent behavior and the dimensionality of the
statistical test. A binary pass/fail verdict discards almost all
information from each trial---the rich execution trace, tool usage
pattern, reasoning depth, and output structure are reduced to a single
bit. By extracting a \emph{behavioral fingerprint} from each
trace---a compact vector capturing the agent's behavioral
signature---we can apply multivariate statistical tests that
extract more information per trial, dramatically reducing the number
of trials needed.

We present three pillars of token-efficient testing: (1)
\emph{behavioral fingerprinting} that enables multivariate regression
detection with higher power per sample; (2) \emph{adaptive budget
optimization} that calibrates the trial count to the agent's actual
behavioral variance rather than worst-case assumptions; and (3)
\emph{trace-first offline analysis} that eliminates live agent
executions entirely for four of six test types by analyzing
previously recorded traces. Together with multi-fidelity proxy
testing and warm-start sequential analysis, these techniques achieve
5--20$\times$ cost reduction while maintaining identical
$(\alpha, \beta)$ statistical guarantees
(\cref{sec:token-efficient}).

\subsection{Contributions}
\label{sec:intro:contributions}

This paper presents \agentassay{}, the first token-efficient framework
for regression testing non-deterministic AI agent workflows, achieving
78--100\% cost reduction at equivalent statistical guarantees, with
86\% behavioral detection power where binary testing has 0\%.
Our contributions are:

\begin{enumerate}[leftmargin=*,label=\textbf{C\arabic*.}]
  \item \textbf{Stochastic Test Semantics} (\cref{sec:stochastic}). We
    define the $(\alpha, \beta, n)$-test triple and a three-valued
    verdict function grounded in confidence intervals and hypothesis
    testing. We prove verdict soundness (\cref{thm:soundness}) and
    regression detection power (\cref{thm:power}). We adapt Wald's
    Sequential Probability Ratio Test (SPRT) for cost-efficient agent
    testing and prove its efficiency advantage (\cref{thm:sprt}).

  \item \textbf{Agent Coverage Metrics} (\cref{sec:coverage}). We
    introduce a five-dimensional coverage tuple
    $\mathcal{C} = (C_{\text{tool}}, C_{\text{path}}, C_{\text{state}},
    C_{\text{boundary}}, C_{\text{model}})$ with formal definitions for
    each dimension and prove coverage monotonicity (\cref{thm:coverage-mono}).

  \item \textbf{Agent Mutation Testing} (\cref{sec:mutation}). We define
    four classes of agent-specific mutation operators
    $\mathcal{M} = \{M_{\text{prompt}}, M_{\text{tool}},
    M_{\text{model}}, M_{\text{context}}\}$, formalize the mutation
    score under stochastic semantics, and prove a mutation adequacy
    theorem (\cref{thm:mutation-adequacy}).

  \item \textbf{Metamorphic Relations} (\cref{sec:metamorphic}). We
    define four families of agent-specific metamorphic relations---
    permutation, perturbation, composition, and oracle---that serve as
    partial test oracles for the oracle problem in agent testing.

  \item \textbf{CI/CD Deployment Gates} (\cref{sec:cicd}). We formalize
    deployment gates as statistical decision procedures with
    user-configurable risk thresholds, integrating stochastic test
    verdicts into continuous deployment pipelines.

  \item \textbf{Contract Integration} (\cref{sec:implementation}). We
    connect \agentassay{} to the \agentassert{}
    framework~\citep{bhardwaj2026abc}, using behavioral contracts as
    formal test oracles and showing that a \PASS{} verdict implies
    $(p, \delta, k)$-satisfaction under specified conditions.

  \item \textbf{Comprehensive Evaluation} (\cref{sec:experiments}). We
    evaluate \agentassay{} across 3 scenarios (e-commerce, customer
    support, code generation), 5 models, and 7{,}605 trials costing
    \$227.

  \item \textbf{Behavioral Fingerprinting}
    (\cref{sec:token:fingerprint}). We define behavioral fingerprints
    that map execution traces to compact vectors on a low-dimensional
    manifold, enabling multivariate regression detection via
    Hotelling's $T^2$ test with provably higher power per sample than
    univariate pass-rate testing
    (\cref{thm:fingerprint-regression}).

  \item \textbf{Adaptive Budget Optimization}
    (\cref{sec:token:budget}). We introduce variance-calibrated budget
    allocation that determines the minimum number of trials for
    $(\alpha, \beta)$-guaranteed testing, achieving 4--7$\times$
    reduction for stable agents (\cref{thm:adaptive-budget}).

  \item \textbf{Trace-First Offline Analysis}
    (\cref{sec:token:trace-first}). We prove that four of six test
    types---coverage, contracts, metamorphic relations, and
    mutation evaluation---can execute at zero additional token cost on
    pre-recorded traces, with formal soundness guarantees
    (\cref{thm:trace-first}).
\end{enumerate}

\subsection{Paper Organization}
\label{sec:intro:organization}

\Cref{sec:related} surveys related work in software testing, statistical
methods, and LLM evaluation. \Cref{sec:stochastic} introduces the formal
stochastic test semantics. \Cref{sec:coverage} defines agent coverage
metrics. \Cref{sec:mutation} presents agent mutation testing.
\Cref{sec:metamorphic} covers metamorphic relations and CI/CD gates.
\Cref{sec:token-efficient} presents the token-efficient testing
framework---behavioral fingerprinting, adaptive budgets, trace-first
analysis, multi-fidelity testing, and warm-start sequential analysis.
\Cref{sec:implementation} describes the \agentassay{} implementation.
\Cref{sec:experiments} reports experimental results.
\Cref{sec:discussion} discusses threats to validity, limitations, and
future work. \Cref{sec:conclusion} concludes. Full proofs are deferred
to \cref{app:proofs}.
